% ---------------------------------------------------------------------------
% Appendix: Multi-turn code instruction following (CodeIF-Bench)
%
% Preamble requirements (add to your main .tex if not already present):
%   \usepackage{booktabs}
%   \usepackage{amsmath}
%   \usepackage{multirow}
%
% Metric macros used throughout this section -- add to your preamble, or keep
% the block below if this file is \input{} before any other use of them:
%   \newcommand{\IA}{\mathrm{IA}}
%   \newcommand{\CA}{\mathrm{CA}}
%   \newcommand{\IFR}{\mathrm{IFR}}
%   \newcommand{\CIF}{\mathrm{CIF}}
%
% BibTeX entry for the citation used below:
%
% @misc{wang2025codeif,
%   title        = {CodeIF-Bench: Evaluating Instruction-Following Capabilities
%                   of Large Language Models in Interactive Code Generation},
%   author       = {Peiding Wang and Li Zhang and Fang Liu and Lin Shi and
%                   Minxiao Li and Bo Shen and An Fu},
%   year         = {2025},
%   eprint       = {2503.22688},
%   archivePrefix= {arXiv},
%   primaryClass = {cs.SE},
%   url          = {https://arxiv.org/abs/2503.22688}
% }
% ---------------------------------------------------------------------------

\providecommand{\IA}{\mathrm{IA}}
\providecommand{\CA}{\mathrm{CA}}
\providecommand{\IFR}{\mathrm{IFR}}
\providecommand{\CIF}{\mathrm{CIF}}

\section{Multi-Turn Code Instruction Following}
\label{app:multiturn-if}

\subsection{Motivation}

Single-turn instruction-following benchmarks measure whether a model complies
with an instruction issued in isolation. They cannot detect the failure mode
that matters for an interactive tutor: satisfying the instruction at turn $t$
by silently undoing the instruction satisfied at turn $t-2$. We therefore
evaluate our candidate models on \textbf{CodeIF-Bench}~\cite{wang2025codeif},
in which every instruction is paired with executable unit tests, so that
instruction compliance is decided by running code rather than by an LLM judge.

\subsection{Benchmark and Task Format}

We use the function-level split of CodeIF-Bench, which is derived from MBPP.
Each task consists of a short natural-language specification (median 101
characters), a set of base assertions, and seven independently verifiable
requirements: \emph{Input--Output Conditions}, \emph{Exception Handling},
\emph{Edge Case Handling}, \emph{Functionality Extension}, \emph{Annotation
Coverage}, \emph{Code Complexity}, and \emph{Code Standard} (PEP~8). Each
requirement carries its own unit tests, which are executed against the current
candidate function.

We exclude \emph{Functionality Extension} from all reported runs. At function
level it deliberately alters the signature (e.g.\ ``add an \texttt{all=True}
parameter''), which invalidates the base assertions against which every other
metric is measured. The reference implementation's dynamic-conversation script
excludes it for the same reason; its static script does not, which is one
source of disagreement between the two.

\paragraph{A caveat on split size.}
The benchmark distributes the function-level split as two files,
\texttt{L\_1\_part\_1.jsonl} (50 tasks) and \texttt{L\_1\_part\_2.jsonl} (20
tasks). Despite the shared \texttt{L\_1} prefix, the second file carries the
\emph{repository-level} schema (\texttt{namespace}, \texttt{project\_path},
\texttt{completion\_path}; its \texttt{tests} field holds \texttt{pytest} node
identifiers rather than runnable assertions). Executing a node-identifier
string as Python yields a bare string literal, which is a valid no-op and exits
successfully: every instruction on those 20 tasks therefore passes vacuously.
We detect and reject these rows, and all results below are computed on the
\textbf{50 genuinely function-level tasks}. We note this explicitly because the
naming invites the error, and because including those rows inflates $\IA$,
$\CA$, $\CIF$ and base retention while deflating $\IFR$.

\subsection{Interaction Protocols}

We evaluate two protocols, both multi-turn:

\begin{description}
  \item[Static.] The requirements are issued one per turn in the benchmark's
    fixed order. No feedback is given and no retries are allowed. This isolates
    instruction retention from error-driven repair.
  \item[Dynamic.] At each round the first requirement that the current code
    fails is issued alone. If the model does not satisfy it, one feedback turn
    is given containing the failing check's error message; if it still fails,
    that requirement is abandoned. This is the reference protocol and is closest
    to a real interactive session.
\end{description}

A third protocol, in which all requirements are stated in a single prompt, is
implemented as a no-conversation control but is not reported here; single-shot
instruction following is covered separately in
Appendix~\ref{app:instruction-following}.

\subsection{Metrics}

Let $S_t$ denote the set of instructions issued up to and including turn $t$
(the base specification counts as one), let $i_t$ be the instruction that turn
$t$ addresses, and let $P_t \subseteq S_{t-1}$ be those instructions the model
had already satisfied at some earlier turn. Writing $\mathbb{1}[n \text{ at } t]$
for the indicator that instruction $n$'s unit tests all pass against turn $t$'s
code:

\begin{align}
  \IA_t  &= \mathbb{1}[i_t \text{ at } t] \\
  \CA_t  &= \frac{1}{|S_t|}\sum_{n \in S_t} \mathbb{1}[n \text{ at } t] \\
  \IFR_t &= \frac{1}{|P_t|}\bigl|\{\, n \in P_t : \neg\,\mathbb{1}[n \text{ at } t] \,\}\bigr| \\
  \CIF   &= \bigl|\{\, n : \mathbb{1}[n \text{ at } T] \,\}\bigr|, \quad T = \text{final turn}
\end{align}

$\IA$ measures compliance with the current instruction; $\CA$ the fraction of
the conversation still satisfied; $\IFR$ the proportion of already-satisfied
instructions that a turn breaks; and $\CIF$ how many instructions survive to the
end. $\CIF$ is computed over all of a task's instructions rather than only those
issued, so that it is comparable across protocols: the dynamic protocol issues
only instructions the current code fails, and an issued-only count would score
it against a harder, self-selected subset.

We additionally report three quantities not present in the reference
implementation:

\begin{itemize}
  \item \textbf{Forgetting by age}: $\IFR$ stratified by $t - u$, where $u$ is
    the turn at which the forgotten instruction was first satisfied. This is the
    context-management curve; a rising curve indicates distance-driven
    forgetting, a flat one indicates that loss depends on instruction type
    rather than on conversational distance.
  \item \textbf{Base retention}: whether the function still passes its original
    MBPP assertions at the final turn. This is the most directly interpretable
    measure of regression and requires no normalisation.
  \item \textbf{Unprompted pass rate}: the proportion of instructions already
    satisfied at turn 1, before being requested. Several requirement types are
    satisfied by any reasonable solution, so $\CA$ and $\CIF$ must be read net
    of this baseline.
\end{itemize}

\paragraph{Comparability.}
$\IA$ should be compared only within a protocol. The dynamic protocol
conditions on failure by construction, so its $\IA$ is measured on a harder
subset than the static protocol's. Cross-protocol comparisons should use
$\CIF$, final $\CA$, or $\IFR$.

\subsection{Experimental Setup}

Each model is served locally with vLLM behind an OpenAI-compatible endpoint
(\texttt{tensor-parallel-size} 1, \texttt{max\_model\_len} 16384, temperature 0,
\texttt{max\_tokens} 2048) on a single NVIDIA A100-80GB. For GPT-OSS-20B we set
\texttt{reasoning\_effort=low}; its responses may place the answer in a
reasoning channel and leave \texttt{content} null, which our client handles by
falling back to \texttt{reasoning\_content}. Candidate code is executed in a
throwaway working directory with a memory cap and a wall-clock timeout; each
unit-test snippet is run as its own subprocess so that a single crash cannot
condemn the remaining snippets of the same requirement.

We reimplemented the harness rather than using the reference scripts, which
hardcode hosted-API credentials and a fixed model name. Four differences affect
the numbers and are stated for transparency: (i) test outcomes are attributed
per snippet rather than per requirement; (ii) a turn producing no runnable code
retains the previous turn's code, whereas the reference implementation treats
raw prose as the candidate and thereby scores a non-answer as total forgetting;
(iii) blank lines are preserved, whereas the reference strips them before
testing, which is self-defeating for the PEP~8 requirement; and (iv) $\CIF$ is
computed as defined above rather than as $\CA_t / t$, a quantity that decays
mechanically with turn index. As a validation, two models were additionally
evaluated with the reference harness on the same 50 tasks; agreement was close
(Gemma-4-31B-it: static $\IA$ 0.773 vs.\ 0.770, $\IFR$ 0.037 vs.\ 0.044;
Qwen2.5-Coder-14B: 0.624 vs.\ 0.655, $\IFR$ 0.037 vs.\ 0.058).

\subsection{Results}

\begin{table}[t]
\centering
\small
\caption{Multi-turn instruction following on CodeIF-Bench (function level, 50
tasks, \emph{Functionality Extension} excluded). Higher is better except
$\IFR$. ``Base ret.'' is the fraction of dialogues whose final function still
passes the original MBPP assertions; ``Unpr.'' is the unprompted pass rate.}
\label{tab:mtif-main}
\begin{tabular}{llcccccc}
\toprule
Model & Protocol & $\IA$ & $\CA$ & $\IFR$ & $\CIF$/7 & Base ret. & Unpr. \\
\midrule
\multirow{2}{*}{Llama-3.1-8B-Instruct}
  & dynamic & 0.269 & 0.315 & 0.343 & 3.04 & 0.360 & 0.453 \\
  & static  & 0.572 & 0.540 & 0.165 & 1.96 & 0.240 & 0.463 \\
\multirow{2}{*}{Qwen2.5-Coder-14B-Instruct}
  & dynamic & 0.318 & 0.536 & 0.022 & 5.02 & 0.760 & 0.480 \\
  & static  & 0.624 & 0.684 & 0.037 & 4.40 & 0.660 & 0.473 \\
\multirow{2}{*}{GPT-OSS-20B}
  & dynamic & 0.569 & 0.686 & 0.048 & 6.20 & 0.820 & 0.567 \\
  & static  & 0.753 & 0.773 & 0.038 & 5.30 & 0.740 & 0.577 \\
\multirow{2}{*}{Qwen3-Coder-30B-A3B-Instruct}
  & dynamic & 0.350 & 0.594 & 0.026 & 5.24 & 0.780 & 0.453 \\
  & static  & 0.661 & 0.755 & 0.031 & 4.68 & 0.700 & 0.456 \\
\multirow{2}{*}{Gemma-4-31B-it}
  & dynamic & \textbf{0.623} & \textbf{0.706} & 0.033 & 6.18 & 0.800 & 0.493 \\
  & static  & \textbf{0.773} & \textbf{0.781} & 0.037 & 5.44 & 0.700 & 0.493 \\
\bottomrule
\end{tabular}
\end{table}

\begin{table}[t]
\centering
\small
\caption{Forgetting by age under the dynamic protocol: $\IFR$ stratified by the
number of turns since the instruction was first satisfied. Bucket sizes for
Llama-3.1-8B are given below the table. The other models' dialogues end
sooner, so their buckets shrink quickly (4--14 observations at ages 6--7);
read those columns as indicative only.}
\label{tab:mtif-age}
\begin{tabular}{lccccccc}
\toprule
Model & 1 & 2 & 3 & 4 & 5 & 6 & 7 \\
\midrule
Llama-3.1-8B-Instruct        & 0.22 & 0.34 & 0.45 & 0.63 & 0.78 & 0.86 & 0.95 \\
Qwen2.5-Coder-14B-Instruct   & 0.04 & 0.07 & 0.07 & 0.06 & 0.08 & 0.07 & 0.17 \\
GPT-OSS-20B                  & 0.03 & 0.08 & 0.16 & 0.21 & 0.19 & 0.00 & 0.20 \\
Qwen3-Coder-30B-A3B-Instruct & 0.03 & 0.03 & 0.06 & 0.07 & 0.09 & 0.21 & 0.75 \\
Gemma-4-31B-it               & 0.02 & 0.06 & 0.09 & 0.11 & 0.09 & 0.11 & 0.25 \\
\midrule
\multicolumn{8}{l}{\footnotesize Llama-3.1-8B bucket sizes: 108, 100, 84, 71, 58, 50, 42} \\
\bottomrule
\end{tabular}
\end{table}

\begin{table}[t]
\centering
\small
\caption{First-attempt $\IA$ by requirement type under the static protocol.}
\label{tab:mtif-type}
\begin{tabular}{lccccccc}
\toprule
Model & Base & I/O & Exc. & Edge & Annot. & Cplx. & PEP~8 \\
\midrule
Llama-3.1-8B-Instruct        & 0.38 & 0.78 & 0.78 & 0.60 & 0.54 & 0.82 & 0.10 \\
Qwen2.5-Coder-14B-Instruct   & 0.48 & 0.80 & 0.86 & 0.73 & 0.42 & 0.92 & 0.16 \\
GPT-OSS-20B                  & 0.54 & 0.86 & 0.98 & 0.75 & 0.68 & 0.96 & 0.50 \\
Qwen3-Coder-30B-A3B-Instruct & 0.54 & 0.90 & 1.00 & 0.77 & 0.42 & 0.96 & 0.04 \\
Gemma-4-31B-it               & 0.54 & 0.84 & 0.98 & 0.81 & 0.84 & 0.94 & 0.46 \\
\bottomrule
\end{tabular}
\end{table}

\subsection{Discussion}

\paragraph{Multi-turn retention is not predicted by single-turn ability.}
% TODO: HumanEval 0.909 is not in this repo's results -- confirm against the
% solving-ability table before submission.
Gemma-4-31B-it and Qwen2.5-Coder-14B-Instruct obtain identical HumanEval
pass@1 (0.909), yet differ by 15 points of static $\IA$ (0.773 vs.\ 0.624) and
by about one instruction in $\CIF$ (5.44 vs.\ 4.40 static). GPT-OSS-20B ranks
second on static $\IA$ and $\CA$ and first on dynamic $\CIF$ and base retention.
Multi-turn instruction following is therefore a distinct axis and is not
recoverable from single-turn coding scores.

\paragraph{$\IFR$ exposes failures that $\IA$ conceals.}
Llama-3.1-8B-Instruct attains a static $\IA$ of 0.572, only 5 points below
Qwen2.5-Coder-14B-Instruct, so on per-turn compliance the two appear broadly
comparable. Their $\IFR$ differs by a factor of $4.5$ (0.165 vs.\ 0.037) and
their base retention by a factor of $2.8$ (0.240 vs.\ 0.660): in 76\% of
dialogues Llama's final function no longer satisfies the specification it was
originally given. It achieves per-turn compliance by destroying prior work, a
pattern no single-turn metric can express.

\paragraph{Only the weakest model forgets with distance.}
$\IFR$ generally rises with instruction age (Table~\ref{tab:mtif-age}), but
the effect differs in kind rather than degree across models. For
Llama-3.1-8B it rises monotonically from 0.22 to 0.95, which points to
distance-driven forgetting: older instructions are progressively lost. The
other four stay at or below 0.21 through age~6. Qwen2.5-Coder-14B is nearly
flat (0.04--0.08), and GPT-OSS-20B and Gemma-4-31B plateau after age~4, so
for these models forgetting depends more on the instruction than on how long
ago it was issued.

\paragraph{Requirement types differ sharply in difficulty.}
PEP~8 compliance is the hardest requirement for every model (0.04--0.50) and
\emph{Annotation Coverage} the most discriminative (0.42--0.84), whereas
\emph{Exception Handling} and \emph{Code Complexity} are near-saturated
(0.78--1.00). The latter two are partly responsible for the high unprompted
pass rates (0.45--0.58): several of their unit tests are satisfied vacuously,
for instance a \texttt{try}/\texttt{except} block with no failure branch, which
passes when no exception is raised. We report the unprompted pass rate so that
$\CA$ and $\CIF$ can be read net of this baseline, and we leave the benchmark's
tests unmodified for comparability.

\subsection{Limitations}

Results cover the function-level split only (50 tasks). The repository-level
splits require the original project checkouts and a working per-project test
environment, and verify requirements through \texttt{pytest} node identifiers
rather than self-contained assertions; they are not evaluated here. Roughly
half of all instructions pass before being requested, so $\CA$ and $\CIF$ should
not be read as pure instruction-following rates. Finally, requirement types are
software-engineering constraints rather than pedagogical ones; transferring this
protocol to tutor-specific instructions is left to future work.
